Badamy problem rozszerzenia zadanej bijekcji pięciu punktów płaszczyzny $\Q^2$ (wierzchołki
kwadratu i punkt $(2,2)$) do automorfizmu wielomianowego o współczynnikach wymiernych
możliwie najmniejszego stopnia; ogólna konstrukcja daje stopień $(m-1)^2=16$.
Najniższy stopień rozwiązania przyjętego przez niezależny weryfikator wynosi 6 (złożenie
odwzorowań trójkątnych stopni 2 i 3).
Korzystając z cytowanego twierdzenia Junga--van der Kulka i udowodnionej multiplikatywności
stopnia, sprowadzamy stopnie $\le5$ do dwóch klas: jednego odwzorowania trójkątnego oraz
złożeń dwóch kwadratowych.
Pierwszą klasę wykluczamy prostym argumentem liniowym, drugą --- certyfikatami
przynależności do ideału, sprawdzonymi niezależnie w Sage i sympy.
Stąd stopień 6 jest minimalny nad $\Q$; dowód jest częściowo wspomagany komputerowo i zależy
od cytowanego twierdzenia.
Nad $\mathbb R$ i $\mathbb C$ minimum wynosi 4 (rozwiązania nad $\Q(\sqrt{94})$), więc różnica
ma charakter arytmetyczny.
Otwarte pozostają m.in. formalna weryfikacja części obliczeniowej oraz opis wszystkich
rozwiązań stopnia 6.
